\documentclass[11pt]{article}
\usepackage{fontspec}
\setmainfont{Times New Roman}
\setsansfont{Arial}
\setmonofont{Consolas}
\usepackage{amsmath,amssymb,mathtools}
\usepackage{geometry}
\usepackage{microtype}
\geometry{a4paper,margin=2.28cm}
\setlength{\parindent}{1.15em}
\setlength{\parskip}{0pt}
\makeatletter
\renewcommand\footnotesize{\@setfontsize\footnotesize{7.5}{8.5}\selectfont}
\makeatother
\providecommand{\mo}{\mathfrak{o}}
\providecommand{\ml}{\mathfrak{l}}
\providecommand{\mc}{\mathfrak{c}}
\providecommand{\mM}{\mathfrak{M}}
\emergencystretch=120em
\allowdisplaybreaks
\sloppy
\hbadness=10000
\vbadness=10000
\hfuzz=2pt
\vfuzz=10pt
\raggedbottom
\begin{document}
% Унит: Папер34 Сецтион18/19 соурце-фиделиты ремедиатион в001. Дирецт Интерславиц Латин ауториты транслатион фром ексистинг сафе Сецтион18 опенинг плус оригинал сцан-витнесс лате Сецтион18 таил and фулл Сецтион19; суперседес цомпрессед Сецтион18 ендинг фор тис боундары.
\section*{34. Хиперкомплексне величины и теорија представјень}
\emph{Продолженье и завршенье: глава III, §§15--19, и глава IV, §§20--26.}

\section*{III. поглавје. Теорија модулов и представјень}

\subsection*{§ 18. Теорија модулов и представјень пополно редуцибилных колц}

Нехај $\mo$ буде пополно редуцибилно и двустранно просто, а $\mM$ конечны $\mo$-модул, за кторы јединичны елемент од $\mo$ јест јединичны оператор. Тогда $\mM$ јест пополно редуцибилны, а просте компоненты сут операторно изоморфне простым левым идеалам $\ml_i$.

Доказ. Нехај
\[
        \mo=\ml_1+\cdots+\ml_n,
        \qquad
        \mM=(\mo m_1,\ldots,\mo m_k),
\]
тогда
\[
        \mM=(\ldots,\ml_i m_k,\ldots).                     \tag{4}
\]
Те $\ml_i m_k$, кторе сут различне од нули, сут изоморфне к $\ml_i$, бо приредјенье $a\mapsto am_k$ јест операторны изоморфизм. Зато $\ml_i m_k$ сут просте. Ако из израза (4) изпустимо те $\ml_i m_k$, кторе уже сут садржане в суми предходных, тогда сума стаје директна.

Последица. Ако к тому $\mM$ јест директно неразложимы, тогда $\mM$ јест просты и изоморфны некому $\ml_i$.

Нехај ныне $\Delta$ буде тело автоморфизмов сеј простој $\mM$, а $\Lambda$ тело автоморфизмов $\ml_i$; тогда на основе операторного изоморфизма меджу $\mM$ и $\ml_i$ имамо изоморфизм колц $\Delta\simeq\Lambda$. По §1 можно сматрјати $\mM$ и $\ml_i$ како бимодулы с $\Delta$ односно $\Lambda$ како правым скаларным подрочјем, при чем ти бимодулы сут такоже модулы представјеньа. Представјеньа посредоване ньими переходет једно в друго чрез изоморфизм колц $\Delta\simeq\Lambda$. Зато за изученье сеј класы представјень можемо ограничити се на представјенье посредовано чрез $\ml_i$ в јего телу автоморфизмов.

Специфично за $\ml_1$ положимо за основу базу из матричных јединиц
\[
        \ml_1=(c_{11},\ldots,c_{n1})
\]
и како реализацију тела автоморфизмов $\Lambda$ возмемо ранеје (§14) с $K$ означено подтело од $\mo$. При том, в сравнаньу с ранејшими размысльеньами, праве идеалы треба заменити левыми, а автоморфизмы одповедајуче писати справа. Ако представимо
\[
        c=\sum_{i,k} c_{ik}\alpha_{ik},
\]
тогда јест
\[
        (cc_{11},\ldots,cc_{n1})=(c_{11},\ldots,c_{n1})(\alpha_{ik}).
\]

\paragraph{Обновјены првотны закльучок §18.}
Ако тогда всакы елемент $a$ од $\mo$ представимо в форме
\[
        a=\sum c_{ik}\alpha_{ik},
\]
тогда $a\mapsto(\alpha_{ik})$ јест представјенье в $K$, посредовано чрез $\ml_1$.

Нехај $\Gamma$ буде подтело од $K$, тако же $K$ јест конечного ранга $t$ взходно к $\Gamma$:
\[
        K=\chi_1\Gamma+\cdots+\chi_t\Gamma .
\]
Тогда $K$ ставаје својим властным модулом представјеньа взходно к $\Gamma$. Елементы $\beta$ од $K$ представјајут се матрицами $B$ в $\Gamma$, кторе се добывајут так:
\[
        \beta\chi_j=\sum_i\chi_i\beta_{ij},\qquad B=(\beta_{ij}).
\]
Ако ныне сматрјати $\ml_1$ како модул представјеньа взходно к подтелу $\Gamma$, находимо:

\emph{Представјенье од $\mo$ в $\Gamma$, посредовано чрез $\ml_1$, добива се тако: в выше поданом представјеньу од $\mo$ в $K$, $a\mapsto(\alpha_{ik})$, елементы $\alpha_{ik}$ заменьајут се матрицами $A_{ik}$, кторе јим одговарјајут в представјеньу од $K$ в $\Gamma$.}

Доказ.
\[
        \ml_1=\sum_i c_{i1}K=\sum_i\sum_j c_{i1}\chi_j\Gamma .
\]
Посредством базы $(\ldots,c_{i1}\chi_j,\ldots)$ елемент $c_{ik}$ представјаје се блоковоју матрицоју, ктора в $(i,k)$-блоку садржи јединичну матрицу $E_t$, а в всих других блоках нулове матрице $0$ реда $t$. Такоже елемент $\alpha$ из $K$ представјаје се блоково-диагоналноју матрицоју с одговарјајучими $t$-редными матрицами $A$. Зато елементу
\[
        a=\sum c_{ik}\alpha_{ik}
\]
одговарја блокова матрица $(A_{ik})$. Преход к либовольному $\mM$ знову даваје изоморфизм колц за представјенье.

Приметка. Тут изучане представјеньа в телу автоморфизмов левых идеалов и в јего конечных подтелах сут, до изоморфизма, једине, кторе могут быти посредоване простыми $\mo$-модулами или левыми идеалами. Бо по ранеје (§2) учиньеној приметце, ако $\mo$-модул $\mM$ сматрјати како бимодул взходно к некојему телу $K$, елементы од $K$ пораждајут хомоморфизмы модулов. Зато тело $K$ необходно мора хомоморфно одображати се в подтело тела автоморфизмов сеј $\mo$-модула. Тој хомоморфизм, понеже јединичны елемент јест јединичны оператор, јест изоморфизм, и цело тело автоморфизмов мора имати конечну степень взходно к одповедајучему подтелу; иначе бы и модул представјеньа имал бесконечны ранг, что не јест можно.

На конец нај буде јешче споменуто, же тут разгледане иредуцибилне представјеньа јешче не сут все, але само те, кторе сут посредоване простыми $\mo$-модулами. Обче сут јешче модулы представјеньа, напр. тело $\Omega$, сматрјано како бимодул взходно к подтелу $\Omega_0$ како левој области и к самому собе како правој области; оне како $\mo$-модулы сут редуцибилне, але како бимодулы иредуцибилне, и зато все-таки водет к иредуцибилным представјеньам.

\subsection*{§ 19. Просте композицијне факторы при модулах и модулах представјеньа}

Нехај $\mo$ буде колцо с условјем максималности и условјем минималности за леве идеалы, а $\mc$ -- максималны нилпотентны идеал. Тогда важи:

\emph{Всакы просты $\mo$-модул $\mM$ или јест анулованы од $\mc$, или јест изоморфны простому левому идеалу $\ml$ из колца без радикала $\mo/\mc$. В последном случају модул јест анулованы всими тыми двустранными идеалами из $\mo/\mc$, кторе не садрже $\ml$, а абсолутна област множительев (§1) јест изоморфна двустранно простому колцу, кторе садржи $\ml$.}

Доказ. Нехај $\mc^\rho=0$. Мора бити $\mc\mM=0$; бо иначе бы было $\mc\mM=\mM$, тогда $\mc^\rho\mM=\mM$, из чего заради $\mc^\rho=0$ следовал бы противречны закльучок $\mM=0$. Зато $\mM$ можно такоже сматрјати како $\mo/\mc$-модул: все елементы једного класа остатков по $\mc$ дају то само при множеньу с елементом из $\mM$.

Како колцо без радикала, $\mo/\mc$ имаје јединичны елемент. Зато $\mM$ јест директна сума модула, кторы јест анулованы од $\mo/\mc$, и модула, за кторы јединичны елемент ставаје јединичны оператор (§17). Понеже $\mM$ јест просты, може наступити само једен из обоју суммандов. Ако јединичны елемент јест јединичны оператор, тогда подобно како в §18 следује изоморфност к простому левому идеалу. Именно, ако $m\ne0$ јест елемент из $\mM$, тогда такоже $(\mo/\mc)m\ne0$; зато $\ml m\ne0$ за најменье једен просты левы идеал $\ml$ из $\mo/\mc$, и зато он мора изчрпати $\mM$. Дальне тврдженьа за таке леве идеалы сут јасне и одмах преносет се на все к ньим изоморфне модулы.

Последок. Ако $\mM$ јест $\mo$-модул, кторы имаје композицијны ред, тогда просте композицијне факторы или сут ануловане од $\mc$, или сут изоморфне простым левым идеалам из $\mo/\mc$.

Ако $P$ јест колцо комутативно повезано с $\mo$, тогда все резултаты оставајут. Треба тогда само вместо модулов сматрјати бимодулы взходно к $P$ справа и $\mo$ слева, кторе задовольајут условју
\[
        (am)\rho=a(m\rho)=(a\rho)m \qquad(\S 15).
\]
Надто јединичны елемент од $P$ такоже имаје быти јединичным оператором. Како подмодулы всегда се разматрајут само таке, кторе допускајут множенье с $P$, како при идеалах. Подмодулы $\mc\mM,\mc^2\mM,\ldots$, сконструоване в наших доказах, задовольајут тому условју.

Особливо, ако $P$ јест тело и $\mM$ јест $P$-модул конечного ранга, тогда $\mM$ јест једночасно модул представјеньа. Представјенье посредовано чрез $\mM$ јест не само рингово-хомоморфно, але такоже операторно-хомоморфно взходно к $P$ (§15). Все представјеньа в $P$ с тој властностју доставльајут се такыми модулами. Иредуцибилне представјеньа принадлежет к простым модулам, и обратно. Зато теоремы сего параграфа можно одмах разумети како теоремы о представјеньах од $\mo$ в $P$: все иредуцибилне представјеньа, кроме нуловых представјень, сут посредоване простыми левыми идеалами од $\mo/\mc$. Ако теды либовольно представјенье по §16 редуцира се посредством композицијного реда, тогда в главној диагонали, кроме нул, наступајут само така представјеньа, кторе сут еквивалентне представјеньам посредованым чрез просте леве идеалы од $\mo/\mc$.

Приметка. Под предположеньами учиньеными о $\mo$ и $P$, за $\mo$ не треба предпоставльати никакого условја конечности. Бо все представјеньа сут једночасно изоморфне представјеньа абсолутној области множительев, а она, како изоморфны праобраз колца из матриц конечного степена, сама ставаје $P$-модул конечного ранга. Понеже по определьеньу за допустиме идеалы сматрјајут се само $P$-модулы, условје максималности и условје минималности за ту абсолутну област множительев сут изполньене.

Та абсолутна област множительев ставаје хиперкомплексны систем взходно к $P$, ако $P$ предполага се како комутативно тело. То предположенье всегда јест изполньено, чим в главној диагонали наступајут не само нулы. Бо тогда $P$, по приметце в §18, јест изоморфно подтелу тела автоморфизмов од $\ml$, кторе при реализацији чрез подтело $K$ од $\mo/\mc$ (§14), заради комутативного повезаньа, мора лежати в центру од $K$.

\end{document}

